% models/mathematical/relativistische_komponenten.tex
% Relativistische Komponenten des integrierten Zeitfeldmodells

\section{Relativistische Komponenten des Zeitfeldes}

\subsection{Grundlage: Raumzeitmetrik}

Die relativistische Schicht des Zeitfeldes basiert auf der lokalen Metrik
\[
ds^2 = g_{\mu\nu}(x) \, dx^\mu dx^\nu.
\]

Die Zeitdilatation ergibt sich aus der Komponente
\[
d\tau = \sqrt{g_{00}} \, dt.
\]

\subsection{Schwarzschild-Metrik}

Für ein sphärisch symmetrisches Gravitationsfeld:
\[
g_{00} = 1 - \frac{2GM}{r}.
\]

Die Zeitdilatation:
\[
\frac{d\tau}{dt} = \sqrt{1 - \frac{2GM}{r}}.
\]

\subsection{Kerr-Metrik}

Für rotierende Massen:
\[
g_{00} = 1 - \frac{2GMr}{\Sigma},
\quad
\Sigma = r^2 + a^2 \cos^2\theta.
\]

Die Rotation koppelt an die Zeitkomponenten über Frame-Dragging.

\subsection{Kosmologische Metrik}

Für großskalige Strukturen:
\[
ds^2 = dt^2 - a(t)^2 \gamma_{ij} dx^i dx^j.
\]

Die kosmologische Zeitrate ergibt sich aus:
\[
d\tau = dt.
\]

\subsection{Relativistische Zeitgradienten}

Die lokalen Zeitgradienten:
\[
\nabla_\mu T^{(r)} = \partial_\mu \sqrt{g_{00}}.
\]

Diese Gradienten koppeln an die planetaren, solaren und galaktischen Felder.

\subsection{Kopplung an das Zeitfeld}

Die relativistische Schicht trägt zum Gesamtzeitfeld bei durch:
\[
\mathcal{T}^{(r)}_{00} = g_{00},
\quad
\mathcal{T}^{(r)}_{0i} = g_{0i},
\quad
\mathcal{T}^{(r)}_{ij} = g_{ij}.
\]

\subsection{Interferenz mit anderen Schichten}

Die Kopplung erfolgt über:
\[
\mathcal{K}\big(\mathcal{T}^{(r)}, \mathcal{T}^{(i)}\big)
= \lambda_{ri} \, g_{\mu\nu} \, \mathcal{T}^{(i)\,\mu\nu}.
\]

\subsection{Ergebnis}

Die relativistische Schicht definiert die lokale Zeitrate und bildet die
metrische Grundlage für alle höheren Zeitfeldkomponenten.
